\chapter{Justificativa}
\label{pre:justificativa}

O principal argumento deste pr\'e-projeto \'e que as redes neurais de grafos t\^em potencial para melhorar significativamente a precis\~ao e a efici\^encia do reconhecimento de padr\~oes educacionais, favorecendo interven\c{c}\~oes precoces e estrat\'egias mais eficazes de acompanhamento. O trabalho pretende explorar de forma detalhada como redes neurais de grafos podem ser aplicadas \`a an\'alise de dados educacionais para identificar padr\~oes no desempenho dos alunos e fatores que influenciam seu progresso acad\^emico.

Al\'em disso, o estudo prop\~oe investigar o impacto pr\'atico dessa abordagem nas institui\c{c}\~oes de ensino e discutir estrat\'egias de interven\c{c}\~ao baseadas nos \textit{insights} gerados por redes neurais de grafos.

O texto do \sigla{TCC}{Trabalho de Conclus\~ao de Curso} seria dividido em se\c{c}\~oes principais, incluindo revis\~ao da literatura sobre redes neurais de grafos e reconhecimento de padr\~oes em dados educacionais, descri\c{c}\~ao detalhada da metodologia, apresenta\c{c}\~ao dos resultados obtidos com conjuntos de dados reais e discuss\~ao das implica\c{c}\~oes pr\'aticas desses resultados. Tamb\'em se previu a an\'alise de limita\c{c}\~oes, desafios e dire\c{c}\~oes futuras de pesquisa.
